% Para 1 — MOTIVATION: Scientific code must move across parallel APIs;
%   practitioners use LLMs; but parallel translation is structurally harder
%   than ordinary translation (thread indexing, sync, memory, host-kernel).
%
% Para 2 — EVALUATION GAP: Existing benchmarks cannot answer the question
%   cleanly. Walk through what each misses:
%   - Sequential benchmarks (HumanEval, SWE-bench): no parallelism
%   - Parallel generation (ParEval): synthesis ≠ translation
%   - Repository-level (ParEval-Repo): 0% — build system is the bottleneck
%   - Kernel-level (LASSI): only 2 directions, coupled to repair mechanism
%   → Pose the central question: "when build infrastructure is fixed, can
%     LLMs translate parallel kernels — and where do they fail?"
%
% Para 3 — OUR APPROACH: Introduce ParBench. Key design choices:
%   - Fix build infrastructure, translate only the kernel
%   - Build-run-verify harness with conjunctive correctness
%   - Survey-grounded corpus (35 repos → 5 suites)
%   - AST-driven augmentation to probe memorization
%
% Para 4 — THREE CONTRIBUTIONS (enumerated):
%   1. Benchmark: 96 specs, 6+4 directions, kernel-centric isolation
%   2. Augmentation engine: 6 transforms × 4 intensity levels
%   3. Empirical characterization: 1,448 records, two models, key findings
%      (a) real capability revealed, (b) build-stage adaptation dominant,
%      (c) direction dependence, (d) augmentation stability

\section{Introduction}
\label{sec:introduction}

Scientific code increasingly needs to move across parallel programming APIs. CUDA implementations written for NVIDIA GPUs must be ported to OpenMP, OpenCL, HIP, SYCL, or other portability layers as hardware and procurement constraints change. Practitioners increasingly use large language models (LLMs) for this work, but parallel translation is not an ordinary code translation problem. Correctness depends on preserving thread-index arithmetic, synchronization, memory-management semantics, and host-kernel interfaces, and more generally, structural concerns that go well beyond surface-level token replacements.

Existing benchmarks do not test the parallel code translation of LLMs. Sequential code benchmarks such as HumanEval~\cite{HumanEval2021} and SWE-bench~\cite{SWEbench2024} contain no parallelism. Parallel code generation benchmarks such as ParEval~\cite{ParEval2024} ask models to synthesize code from natural language descriptions and do not test translation of existing parallel programs across APIs. Repository-level translation, as in ParEval-Repo~\cite{ParEvalRepo2025}, evaluates a repository-level translation but makes build-system and integration generation a necessary task, reporting 0\% pass@$k$ on applications larger than 133~SLoC. LASSI~\cite{LASSI2024} takes a step towards kernel-level evaluation with an agentic self-correction pipeline on 10 HeCBench kernels, but covers only two translation directions and couples evaluation to the repair mechanism. This leaves two questions open: can models translate the parallel computational kernel itself across a broad set of directions when build infrastructure is fixed? And do the results reflect genuine  reasoning of parallel code semantics, or are LLMs simply memorizing publicly-available benchmark code?

To address these gaps, we present \textbf{ParBench}, a kernel-centric benchmark framework for evaluating LLM-based parallel code translation. ParBench provides a fixed build and verification infrastructure, asks the model to translate only the selected kernel files, and judges the output with a build-run-verify harness. The benchmark corpus is not ad-hoc: it is derived from a survey of 35 open-source HPC benchmark repositories and narrowed to suites with same kernels written in multiple APIs and automated correctness checks. ParBench's design addresses several dimensions that prior work leaves open:

\textbf{Survey-grounded benchmark selection.} Prior benchmarks do not justify \emph{which} parallel APIs or kernels to evaluate. ParBench's survey analyzed kernel-level translation opportunities across the 35 repositories and revealed that kernel-level counting exposes far more translation material than repository-level counting. For example, just six repositories yielded 472 CUDA--OpenMP kernel pairs (Section~\ref{sec:suite-selection}; Appendix Figure~\ref{fig:repo-vs-kernel}). From this analysis, five suites were selected for their multi-API kernel equivalence and self-checking correctness infrastructure: Rodinia~\cite{Rodinia2009}, HeCBench~\cite{HeCBench2023}, XSBench~\cite{XSBench2014}, RSBench, and mixbench~\cite{mixbench2017}, spanning CUDA, OpenMP, and OpenCL with additional OpenMP-target case studies.

\textbf{Multi-file coordination.} A substantial fraction of ParBench specifications require translating two or more source files, testing whether models can coordinate changes across file boundaries, a challenge distinct from within-file API mapping. Single-file translations pass at substantially higher rates, suggesting that file-boundary coordination adds a separate challenge (Section~\ref{sec:eval-corpus}).

\textbf{Training data contamination.} HPC suites such as Rodinia and HeCBench are widely available across hundreds of GitHub repositories, so translation success on these benchmarks may reflect memorization by LLMs rather than their robust reasoning of parallel code semantics. To address this issue, ParBench introduces an AST-driven augmentation engine that produces semantically-equivalent but syntactically novel source variants of varying degrees, enabling controlled evaluation of whether models reason about parallel code semantics or rely on syntactic similarity (Section~\ref{sec:augmentation-engine}; results in Section~\ref{sec:augmentation-robustness}).

This paper makes three contributions:

\begin{enumerate}
\item \textbf{Parallel code translation benchmark.} We introduce ParBench, a benchmark for LLM-based parallel code translation that provides executable kernel specifications spanning five HPC benchmark suites with six standard translation directions among CUDA, OpenMP, and OpenCL, plus four OpenMP-target case-study directions. By fixing the surrounding build infrastructure as context, ParBench enables direct build-run-verify evaluation of kernel translation without the build-system-generation confound identified by repository-level work.

\item \textbf{AST-driven augmentation engine.} We develop an augmentation engine with six semantics-preserving source-level transformations applied at four intensity levels (L1--L4). This engine tests whether models preserve parallel semantics beyond surface similarity by producing structurally varied but functionally equivalent source code.

% src: quantitative_findings_qwen.json + quantitative_findings_gpt54.json > canonical > pass_at_k
\item \textbf{Empirical characterization.} Across two models and 1,448 valid translation records, we report the first multi-suite, multi-direction evaluation of LLM parallel code translation under stochastic sampling with pass@$k$ analysis. \qwenshort{} achieves pass@1\,=\,23.9\% and \gptnew{} achieves pass@1\,=\,62.7\% on 142 unique no-augmentation (L0) tasks. Key findings include build-stage adaptation as the main bottleneck, strong direction dependence, broad stability under augmentation, and limited rescue from repeated sampling.
\end{enumerate}

{\color{blue}
Scientific code increasingly needs to move across parallel programming APIs as hardware platforms, software ecosystems, and portability requirements evolve. For example, CUDA implementations may need to be ported to OpenMP, OpenCL, or other portability layers when applications target new accelerators or seek broader hardware support. Large language models (LLMs) are increasingly being explored for this task, but parallel code translation is fundamentally different from ordinary code translation. Its correctness depends not only on preserving program intent, but also on maintaining thread-index arithmetic, synchronization semantics, memory-management patterns, and host-kernel interfaces—structural properties that go far beyond surface-level token replacement.

Can current LLMs actually perform this translation? Existing evaluations leave this question unresolved because their benchmark datasets either lack parallel semantics or focus on generation rather than translation. Widely adopted code-generation benchmarks such as HumanEval and SWE-bench~\cite{HumanEval2021, SWEbench2024} are predominantly sequential and therefore do not evaluate parallel programming semantics. Parallel code-generation benchmarks such as ParEval~\cite{ParEval2024} primarily evaluate synthesis from natural-language descriptions, rather than translation of existing parallel implementations across APIs. This leaves two questions open: can models translate computational kernels across a broad set of API directions when the build infrastructure is fixed, and do successful translations reflect genuine parallel reasoning rather than memorization of widely circulated benchmark code?

We present \textbf{ParBench}, a benchmark framework that answers these questions by isolating parallel kernel translation from repository reconstruction. ParBench fixes the build and verification infrastructure, asks models to translate only computational kernel files, and evaluates each generated kernel with a build-run-verify harness using conjunctive correctness checks. The corpus is survey-grounded: we review 35 open-source HPC repositories and select five benchmark suites that provide equivalent multi-API kernels and automated correctness checks. To test whether successful translations reflect genuine parallel reasoning rather than memorization, ParBench includes an AST-driven augmentation engine that generates semantically equivalent but syntactically novel source variants.
This paper makes three contributions:

\begin{enumerate}
\item \textbf{Parallel code translation benchmark.} We introduce \textbf{ParBench}, a benchmark for LLM-based parallel code translation that provides executable kernel specifications spanning five HPC benchmark suites, six standard translation directions among CUDA, OpenMP, and OpenCL, and four OpenMP-target case-study directions. By fixing the surrounding build infrastructure as context, ParBench enables direct build-run-verify evaluation of kernel translation without the build-system-generation confound identified by repository-level work.

\item \textbf{AST-driven augmentation engine.} We develop an augmentation engine with six semantics-preserving source-level transformations applied at four intensity levels (L1--L4). The engine produces structurally varied but functionally equivalent source variants, allowing us to test whether models preserve parallel semantics beyond surface similarity to widely circulated benchmark code.

\item \textbf{Empirical characterization.} Across two models and 1,448 valid translation records, we report the first multi-suite, multi-direction evaluation of LLM-based parallel code translation under stochastic sampling with pass@$k$ analysis. \qwenshort{} achieves pass@1\,=\,23.9\% and \gptnew{} achieves pass@1\,=\,62.7\% on 142 unique no-augmentation (L0) tasks. Key findings include build-stage API adaptation as the main bottleneck, strong direction and file-boundary effects, broad stability under augmentation, and limited gains from repeated sampling.
\end{enumerate}

}